% GENERATED FILE. Edit canonical lessons, then run npm run generate:books.
% canonical-source-hash: fnv1a64:b860d4c14fdb99b5
% canonical-lessons: IT-C458-raccomandata, IT-C458-spedizione, IT-C458-mittente, IT-C458-affrancare, IT-C458-recapitare

\chapter{Alla posta}
\label{ch:it-alla-posta}
\hlchaptermodality{458}

\begin{chapteropening}
\textbf{By the end of this chapter:} \emph{I can send things by post: a registered letter, the shipping, the sender, the postage and delivery.}

Complete the last lesson of chapter 458: I can send things by post: a registered letter, the shipping, the sender, the postage and delivery.
\end{chapteropening}

\section[la raccomandata]{la raccomandata — the registered letter}

\label{lesson:IT-C458-raccomandata}



\begin{quote}
Before the new one: say the Italian for the clearance sale, then the Italian for to haggle.
\end{quote}

\subsection*{You'll want to know: la raccomandata}
\textbf{la raccomandata} — \textquotedblleft{}the registered letter\textquotedblright{}.

\begin{quote}
Ho spedito una raccomandata con ricevuta di ritorno.
\end{quote}

\begin{quote}
\emph{I sent a registered letter with a return receipt.}
\end{quote}

\emph{Raccomandare} — to recommend, and to register a letter.

\begin{cousinweb}
\emph{raccomandata} is built on \emph{raccomandare} 'to entrust' (EtymDB), from \emph{ri-} and \emph{comandare}, which goes back to Latin \emph{commendāre} 'to entrust', built on \emph{mandāre}; the words are inherited.
\end{cousinweb}

\begin{grammarlens}[title={What you've built}]
One more word for telling, explaining and describing at B1.
\end{grammarlens}

\subsection*{Your turn}
\begin{itemize}
\raggedright
  \item \emph{Say it:} \emph{la raccomandata}
  \item \emph{Say it:} \emph{la raccomandata}, once more
  \item \emph{Say it:} a sentence of your own with \emph{la raccomandata}
  \item \emph{Recall:} say the Italian for the clearance sale, then the Italian for to haggle, then say \emph{la raccomandata} again
\end{itemize}

\begin{tcolorbox}[breakable,colback=teal!4,colframe=teal!35!black,title={Before you move on}]
What does la raccomandata mean? (\textquotedblleft{}The registered letter\textquotedblright{}.) Say it once more.
\end{tcolorbox}
\section[la spedizione]{la spedizione — the shipping, the dispatch; the expedition}

\label{lesson:IT-C458-spedizione}



\begin{quote}
Before the new one: say the Italian for to haggle, then the Italian for the registered letter.
\end{quote}

\subsection*{You'll want to know: la spedizione}
\textbf{la spedizione} — \textquotedblleft{}the shipping, the dispatch; the expedition\textquotedblright{}.

\begin{quote}
La spedizione costa cinque euro.
\end{quote}

\begin{quote}
\emph{Shipping costs five euros.}
\end{quote}

Also an expedition to a mountain or to the poles.

\begin{cousinweb}
\emph{spedizione} is a learned borrowing from Latin \emph{expedītiō} 'a setting out' (EtymDB), from \emph{expedīre} 'to free the feet, to get ready', built on \emph{pēs} 'foot'.
\end{cousinweb}

\begin{grammarlens}[title={What you've built}]
One more word for telling, explaining and describing at B1.
\end{grammarlens}

\subsection*{Your turn}
\begin{itemize}
\raggedright
  \item \emph{Say it:} \emph{la spedizione}
  \item \emph{Say it:} \emph{la spedizione}, once more
  \item \emph{Say it:} a sentence of your own with \emph{la spedizione}
  \item \emph{Recall:} say the Italian for to haggle, then the Italian for the registered letter, then say \emph{la spedizione} again
\end{itemize}

\begin{tcolorbox}[breakable,colback=teal!4,colframe=teal!35!black,title={Before you move on}]
What does la spedizione mean? (\textquotedblleft{}The shipping, the dispatch; the expedition\textquotedblright{}.) Say it once more.
\end{tcolorbox}
\section[il mittente]{il mittente — the sender}

\label{lesson:IT-C458-mittente}



\begin{quote}
Before the new one: say the Italian for the registered letter, then the Italian for the shipping.
\end{quote}

\subsection*{You'll want to know: il mittente}
\textbf{il mittente} — \textquotedblleft{}the sender\textquotedblright{}.

\begin{quote}
Il mittente di questa lettera è mia zia.
\end{quote}

\begin{quote}
\emph{The sender of this letter is my aunt.}
\end{quote}

The person you send to is \emph{il destinatario}.

\begin{cousinweb}
\emph{mittente} is a learned borrowing from Latin \emph{mittēns} 'sending', from \emph{mittere} 'to send' (Lewis \& Short).
\end{cousinweb}

\begin{grammarlens}[title={What you've built}]
One more word for telling, explaining and describing at B1.
\end{grammarlens}

\subsection*{Your turn}
\begin{itemize}
\raggedright
  \item \emph{Say it:} \emph{il mittente}
  \item \emph{Say it:} \emph{il mittente}, once more
  \item \emph{Say it:} a sentence of your own with \emph{il mittente}
  \item \emph{Recall:} say the Italian for the registered letter, then the Italian for the shipping, then say \emph{il mittente} again
\end{itemize}

\begin{tcolorbox}[breakable,colback=teal!4,colframe=teal!35!black,title={Before you move on}]
What does il mittente mean? (\textquotedblleft{}The sender\textquotedblright{}.) Say it once more.
\end{tcolorbox}
\section[affrancare]{affrancare — to put postage on; to set free}

\label{lesson:IT-C458-affrancare}



\begin{quote}
Before the new one: say the Italian for the shipping, then the Italian for the sender.
\end{quote}

\subsection*{You'll want to know: affrancare}
\textbf{affrancare} — \textquotedblleft{}to put postage on; to set free\textquotedblright{}.

\begin{quote}
Devo affrancare questa lettera per l'estero.
\end{quote}

\begin{quote}
\emph{I need to put postage on this letter going abroad.}
\end{quote}

Its first meaning is to set free: \emph{franco} once meant free.

\begin{cousinweb}
\emph{affrancare} is built from \emph{a-} and \emph{franco} 'free', which was borrowed from Frankish, the Germanic language of the Franks: a free man was a Frank.
\end{cousinweb}

\begin{grammarlens}[title={What you've built}]
One more word for telling, explaining and describing at B1.
\end{grammarlens}

\subsection*{Your turn}
\begin{itemize}
\raggedright
  \item \emph{Say it:} \emph{affrancare}
  \item \emph{Say it:} \emph{affrancare}, once more
  \item \emph{Say it:} a sentence of your own with \emph{affrancare}
  \item \emph{Recall:} say the Italian for the shipping, then the Italian for the sender, then say \emph{affrancare} again
\end{itemize}

\begin{tcolorbox}[breakable,colback=teal!4,colframe=teal!35!black,title={Before you move on}]
What does affrancare mean? (\textquotedblleft{}To put postage on; to set free\textquotedblright{}.) Say it once more.
\end{tcolorbox}
\section[recapitare]{recapitare — to deliver (mail)}

\label{lesson:IT-C458-recapitare}



\begin{quote}
Before the new one: say the Italian for the sender, then the Italian for to put postage on.
\end{quote}

\subsection*{You'll want to know: recapitare}
\textbf{recapitare} — \textquotedblleft{}to deliver (mail)\textquotedblright{}.

\begin{quote}
Il postino ha recapitato il pacco al vicino.
\end{quote}

\begin{quote}
\emph{The postman delivered the parcel to the neighbour.}
\end{quote}

\emph{Il recapito} — the address, the contact details.

\begin{cousinweb}
\emph{recapitare} is built on \emph{recapito} 'address, delivery'; the deeper origin is not settled.
\end{cousinweb}

\begin{grammarlens}[title={What you've built}]
One more word for telling, explaining and describing at B1.
\end{grammarlens}

\subsection*{Your turn}
\begin{itemize}
\raggedright
  \item \emph{Say it:} \emph{recapitare}
  \item \emph{Say it:} \emph{recapitare}, once more
  \item \emph{Say it:} a sentence of your own with \emph{recapitare}
  \item \emph{Recall:} say the Italian for the sender, then the Italian for to put postage on, then say \emph{recapitare} again
\end{itemize}

\begin{tcolorbox}[breakable,colback=teal!4,colframe=teal!35!black,title={Before you move on}]
What does recapitare mean? (\textquotedblleft{}To deliver (mail)\textquotedblright{}.) Say it once more.
\end{tcolorbox}
